\documentclass{article}
\usepackage{amsmath,amssymb,amsthm}
\usepackage{algorithm}
\usepackage[noend]{algpseudocode}
\usepackage{hyperref}
\usepackage{enumitem}
\newtheorem{theorem}{Theorem}
\newtheorem{lemma}{Lemma}
\newtheorem{assumption}{Assumption}
\newtheorem{remark}{Remark}
\begin{document}

\title{AMSGrad: Algorithmic Details and Full Convergence Proof for Non‑Convex Smooth Objectives}
\author{}
\date{}
\maketitle

%%%%%%%%%%%%%%%%%%%%%%%%%%%%%%%%%%%%%%%%%%%%%%%%%%%%%%%%%%%%
\section*{Algorithm Name}
\textbf{AMSGrad} (Adaptive Moment Estimation with a Monotonically Non‑Increasing Second Moment)

%%%%%%%%%%%%%%%%%%%%%%%%%%%%%%%%%%%%%%%%%%%%%%%%%%%%%%%%%%%%
\section*{Mathematical Formulation}
Let $\{z_t\}_{t\ge 1}$ denote an i.i.d. data stream and let  
$g_t:=\nabla f(w_t;z_t)$ be a stochastic gradient of the (possibly non‑convex) loss $f(\cdot)$.  
Hyper‑parameters are:
\[
\eta>0\quad\text{(base learning rate)},\qquad 
\beta_1\in[0,1),\ \beta_2\in[0,1),\qquad
\epsilon>0\quad\text{(numerical stability)} .
\]

The iterates are defined for $t=1,2,\dots$ by
\begin{subequations}\label{eq:amsgrad}
\begin{align}
m_t &:= \beta_1 m_{t-1} + (1-\beta_1) g_t, \tag{1}\\[4pt]
v_t &:= \beta_2 v_{t-1} + (1-\beta_2) g_t\odot g_t, \tag{2}\\[4pt]
\widehat v_t &:= \max\!\bigl(\widehat v_{t-1},\,v_t\bigr) \quad\text{(element‑wise max)},\tag{3}\\[4pt]
\widehat m_t &:= \frac{m_t}{1-\beta_1^{\,t}},\qquad
\widehat v_t^{\;{\rm corr}}:=\frac{\widehat v_t}{1-\beta_2^{\,t}},\tag{4}\\[4pt]
w_{t+1} &:= w_t - \eta\,
\frac{\widehat m_t}{\sqrt{\widehat v_t^{\;{\rm corr}}}+\epsilon}. \tag{5}
\end{align}
\end{subequations}
All operations are component‑wise; $\odot$ denotes element‑wise product.

%%%%%%%%%%%%%%%%%%%%%%%%%%%%%%%%%%%%%%%%%%%%%%%%%%%%%%%%%%%%
\section*{Pseudocode}
\begin{algorithm}[H]
\caption{AMSGrad}
\begin{algorithmic}[1]
\Require Initial point $w_0\in\mathbb{R}^d$, step size $\eta>0$, decay rates $\beta_1,\beta_2\in[0,1)$, $\epsilon>0$.
\State $m_0\gets 0,\; v_0\gets 0,\; \widehat v_0\gets 0$
\For{$t=1,2,\dots,T$}
    \State Sample mini‑batch (or a single example) and compute stochastic gradient $g_t=\nabla f(w_t;z_t)$
    \State $m_t \gets \beta_1 m_{t-1} + (1-\beta_1) g_t$ \Comment{first moment}
    \State $v_t \gets \beta_2 v_{t-1} + (1-\beta_2) g_t\odot g_t$ \Comment{second moment}
    \State $\widehat v_t \gets \max(\widehat v_{t-1},\,v_t)$ \Comment{AMSGrad correction}
    \State $\widehat m_t \gets m_t/(1-\beta_1^{\,t})$ \Comment{bias‑corrected first moment}
    \State $\widehat v_t^{\;{\rm corr}} \gets \widehat v_t/(1-\beta_2^{\,t})$ \Comment{bias‑corrected second moment}
    \State $w_{t+1} \gets w_t - \eta\,\widehat m_t\;/\;\bigl(\sqrt{\widehat v_t^{\;{\rm corr}}}+\epsilon\bigr)$
\EndFor
\State \Return $w_{T+1}$
\end{algorithmic}
\end{algorithm}

%%%%%%%%%%%%%%%%%%%%%%%%%%%%%%%%%%%%%%%%%%%%%%%%%%%%%%%%%%%%
\section*{Dry Run Example}
Consider the simple one‑dimensional quadratic
\[
f(w) = (w-3)^2,\qquad \nabla f(w) = 2(w-3).
\]
We use a deterministic gradient ($g_t=\nabla f(w_t)$), base step size $\eta=0.1$, $\beta_1=0.9$, $\beta_2=0.999$, $\epsilon=10^{-8}$, and start from $w_0=0$.

\begin{center}
\begin{tabular}{c|c|c|c|c|c}
$t$ & $w_t$ & $g_t$ & $m_t$ & $v_t$ & $w_{t+1}$\\ \hline
0 & $0$ & – & $0$ & $0$ & – \\
1 & $0$ & $-6$ & $0.1\cdot0+0.9(-6)=-5.4$ & $0.001\cdot0+0.999\cdot36=35.964$ & 
$0-0.1\cdot\frac{-5.4}{\sqrt{35.964}+10^{-8}}\approx0.0900$\\
2 & $0.0900$ & $2(0.0900-3)=-5.82$ &
$m_2=0.9(-5.4)+0.1(-5.82)=-5.562$ &
$v_2=0.999\cdot35.964+0.001\cdot33.8724\approx35.928$ &
$w_3\approx0.0900-0.1\frac{-5.562}{\sqrt{35.928}}\approx0.1763$\\
3 & $0.1763$ & $2(0.1763-3)=-5.647$ &
$m_3=0.9(-5.562)+0.1(-5.647)=-5.553$ &
$v_3=0.999\cdot35.928+0.001\cdot31.872\approx35.896$ &
$w_4\approx0.1763-0.1\frac{-5.553}{\sqrt{35.896}}\approx0.2595$
\end{tabular}
\end{center}
Objective values: $f(w_0)=9$, $f(w_1)\approx8.28$, $f(w_2)\approx7.66$, $f(w_3)\approx7.12$.  The iterates move toward the minimiser $w^\star=3$.

%%%%%%%%%%%%%%%%%%%%%%%%%%%%%%%%%%%%%%%%%%%%%%%%%%%%%%%%%%%%
\section*{Assumptions}
\begin{assumption}[Smoothness]\label{as:smooth}
The stochastic loss $f(\cdot;z)$ is $L$‑smooth for every $z$, i.e.
\[
\|\nabla f(w;z)-\nabla f(u;z)\|\le L\|w-u\|,\qquad\forall w,u\in\mathbb{R}^d.
\]
Consequently, the expected objective $F(w):=\mathbb{E}_z[f(w;z)]$ is also $L$‑smooth.
\end{assumption}
\begin{assumption}[Bounded stochastic gradients]\label{as:bounded}
There exists $G>0$ such that
\[
\|g_t\|_\infty\le G\quad\text{almost surely for all }t.
\]
\end{assumption}
\begin{assumption}[Unbiasedness and bounded variance]\label{as:unbiased}
\[
\mathbb{E}\!\left[g_t\mid\mathcal{F}_{t-1}\right]=\nabla F(w_t),\qquad
\mathbb{E}\!\left[\|g_t-\nabla F(w_t)\|^2\mid\mathcal{F}_{t-1}\right]\le\sigma^2,
\]
where $\mathcal{F}_{t-1}$ denotes the sigma‑algebra generated by $\{w_s,g_s\}_{s\le t-1}$.
\end{assumption}
\begin{assumption}[Hyper‑parameter constraints]\label{as:params}
\[
\beta_1\in[0,1),\qquad \beta_2\in[0,1),\qquad
\eta_t:=\frac{\eta}{\sqrt{t}}\ \text{with }\eta>0.
\]
Moreover we require $\beta_1<\sqrt{\beta_2}$ (standard condition ensuring the “effective” decay of the first moment).
\end{assumption}

%%%%%%%%%%%%%%%%%%%%%%%%%%%%%%%%%%%%%%%%%%%%%%%%%%%%%%%%%%%%
\section*{Main Convergence Theorem}
\begin{theorem}[Non‑convex convergence of AMSGrad]\label{thm:main}
Assume \ref{as:smooth}–\ref{as:params}.  Let $\{w_t\}$ be generated by \eqref{eq:amsgrad} with step sizes $\eta_t=\eta/\sqrt{t}$.  Define
\[
\overline{w}_T:=\frac{1}{T}\sum_{t=1}^{T}w_t .
\]
Then for any $T\ge 1$,
\[
\frac{1}{T}\sum_{t=1}^{T}\mathbb{E}\!\bigl[\|\nabla F(w_t)\|^2\bigr]
\;\le\;
\underbrace{\frac{2L\bigl(F(w_1)-F^\star\bigr)}{\eta\sqrt{T}}}_{\text{optimization term}}
\;+\;
\underbrace{\frac{C\,\eta\, G\,(1+\log T)}{\sqrt{T}}}_{\text{variance term}},
\]
where $F^\star:=\inf_wF(w)$ and the constant
\[
C:=\frac{2\bigl(1-\beta_1\bigr)}{(1-\beta_2)(1-\beta_1/\sqrt{\beta_2})}
\]
depends only on $\beta_1,\beta_2$.  Consequently,
\[
\min_{1\le t\le T}\mathbb{E}\!\bigl[\|\nabla F(w_t)\|^2\bigr]
= O\!\left(\frac{\log T}{\sqrt{T}}\right).
\]
\end{theorem}

\begin{remark}
When the stochastic gradient noise vanishes ($\sigma=0$) the second term disappears and the bound reduces to $O(1/\sqrt{T})$, matching the deterministic rate for smooth non‑convex optimisation.
\end{remark}

%%%%%%%%%%%%%%%%%%%%%%%%%%%%%%%%%%%%%%%%%%%%%%%%%%%%%%%%%%%%
\section*{Theoretical Properties}
\begin{itemize}[leftmargin=*]
\item \textbf{Stability.} The monotone max in \eqref{eq:amsgrad} guarantees that each denominator $\sqrt{\widehat v_t^{\;{\rm corr}}}$ is never smaller than its past values, preventing the effective step size from exploding.
\item \textbf{Hyper‑parameter sensitivity.} The constant $C$ in Theorem~\ref{thm:main} shows explicitly how $\beta_1$ and $\beta_2$ influence the bound.  Larger $\beta_2$ (slower decay of the second moment) reduces $C$ but also makes the denominator larger, yielding a more conservative step.
\item \textbf{Comparison to Adam.} Adam uses $v_t$ directly (no max).  Under the same assumptions Adam may fail to converge for non‑convex problems (Reddi et al., 2018).  The AMSGrad correction restores a provable $O(\log T/\sqrt{T})$ guarantee.
\item \textbf{Special cases.} If $\beta_1=0$, AMSGrad reduces to a variant of RMSProp with monotone second moment; the bound simplifies with $C=2/(1-\beta_2)$.
\end{itemize}

%%%%%%%%%%%%%%%%%%%%%%%%%%%%%%%%%%%%%%%%%%%%%%%%%%%%%%%%%%%%
\section*{Preliminaries (Definitions and Lemmas)}
\begin{lemma}[Smoothness inequality]\label{lem:smooth}
For any $L$‑smooth $F$ and any $w,u$,
\[
F(u) \le F(w) + \langle\nabla F(w),u-w\rangle + \frac{L}{2}\|u-w\|^2 .
\]
\end{lemma}
\begin{proof}
Standard consequence of $L$‑smoothness; see e.g. Nesterov (2004). \qedhere
\end{proof}

\begin{lemma}[Bounding the adaptive denominator]\label{lem:denom}
For all $t$ and each coordinate $i$,
\[
\sqrt{\widehat v_{t,i}^{\;{\rm corr}}} \ge \sqrt{(1-\beta_2)G^2}\;,
\]
hence
\[
\frac{1}{\sqrt{\widehat v_{t,i}^{\;{\rm corr}}}} \le \frac{1}{\sqrt{(1-\beta_2)}\,G}.
\]
\end{lemma}
\begin{proof}
From \eqref{eq:amsgrad} and the boundedness assumption $\|g_t\|_\infty\le G$ we have
\[
v_{t,i} = \beta_2 v_{t-1,i} + (1-\beta_2)g_{t,i}^2 \ge (1-\beta_2)g_{t,i}^2\ge (1-\beta_2)G^2.
\]
Since $\widehat v_t\ge v_t$ element‑wise, the claim follows after bias correction. \qedhere
\end{proof}

\begin{lemma}[Bias‑correction lower bound]\label{lem:bias}
For any $t\ge 1$,
\[
\frac{1}{1-\beta_2^{\,t}} \le \frac{1}{1-\beta_2}.
\]
\end{lemma}
\begin{proof}
Because $0\le\beta_2<1$, the sequence $\beta_2^{\,t}$ is decreasing, hence $1-\beta_2^{\,t}\ge 1-\beta_2$. \qedhere
\end{proof}

%%%%%%%%%%%%%%%%%%%%%%%%%%%%%%%%%%%%%%%%%%%%%%%%%%%%%%%%%%%%
\section*{Main Proof (Step‑by‑Step Derivation)}
We prove Theorem~\ref{thm:main}.  Throughout the proof, expectations are taken conditionally on $\mathcal{F}_{t-1}$ unless otherwise noted.

\bigskip
\noindent\textbf{[Step 1] Apply smoothness.}
Using Lemma~\ref{lem:smooth} with $u=w_{t+1}$,
\[
F(w_{t+1}) \le F(w_t) + \big\langle\nabla F(w_t), w_{t+1}-w_t\big\rangle
               + \frac{L}{2}\|w_{t+1}-w_t\|^2 . \tag{6}
\]

\noindent\textbf{[Step 2] Substitute the AMSGrad update.}
From \eqref{eq:amsgrad} (5),
\[
w_{t+1}-w_t = -\eta_t\,\frac{\widehat m_t}{\sqrt{\widehat v_t^{\;{\rm corr}}}+\epsilon}.
\]
Insert into (6):
\[
F(w_{t+1}) \le F(w_t) -\eta_t\,
\Big\langle\nabla F(w_t),\frac{\widehat m_t}{\sqrt{\widehat v_t^{\;{\rm corr}}}+\epsilon}\Big\rangle
+ \frac{L\eta_t^{2}}{2}\Big\|\frac{\widehat m_t}{\sqrt{\widehat v_t^{\;{\rm corr}}}+\epsilon}\Big\|^{2}. \tag{7}
\]

\noindent\textbf{[Step 3] Take conditional expectation.}
Using unbiasedness (Assumption~\ref{as:unbiased}) and the fact that $\widehat m_t$ is $\mathcal{F}_{t-1}$‑measurable,
\[
\mathbb{E}\!\bigl[F(w_{t+1})\mid\mathcal{F}_{t-1}\bigr]
\le F(w_t) -\eta_t
\Big\langle\nabla F(w_t),\frac{m_t}{1-\beta_1^{t}}\frac{1}{\sqrt{\widehat v_t^{\;{\rm corr}}}+\epsilon}\Big\rangle
+ \frac{L\eta_t^{2}}{2}\,\mathbb{E}\!\Bigl[\big\|\frac{\widehat m_t}{\sqrt{\widehat v_t^{\;{\rm corr}}}+\epsilon}\big\|^{2}\Big|\mathcal{F}_{t-1}\Bigr]. \tag{8}
\]

\noindent\textbf{[Step 4] Relate $m_t$ to the true gradient.}
From the definition of $m_t$,
\[
m_t = (1-\beta_1)g_t + \beta_1 m_{t-1}.
\]
Iterating and using $m_0=0$ yields
\[
m_t = (1-\beta_1)\sum_{i=1}^{t}\beta_1^{t-i}g_i .
\]
Thus,
\[
\frac{m_t}{1-\beta_1^{t}} = \sum_{i=1}^{t}\alpha_{t,i} g_i,\qquad
\alpha_{t,i}:=\frac{(1-\beta_1)\beta_1^{t-i}}{1-\beta_1^{t}}.
\]
Note that $\alpha_{t,i}\ge0$ and $\sum_{i=1}^{t}\alpha_{t,i}=1$.

\noindent\textbf{[Step 5] Decompose the inner product.}
Using linearity of expectation and the fact that $\mathbb{E}[g_i\mid\mathcal{F}_{t-1}]=\nabla F(w_i)$,
\[
\mathbb{E}\!\Bigl[\Big\langle\nabla F(w_t),\frac{m_t}{1-\beta_1^{t}}\frac{1}{\sqrt{\widehat v_t^{\;{\rm corr}}}+\epsilon}\Big\rangle\Big|\mathcal{F}_{t-1}\Bigr]
= \Big\langle\nabla F(w_t),\frac{1}{\sqrt{\widehat v_t^{\;{\rm corr}}}+\epsilon}\sum_{i=1}^{t}\alpha_{t,i}\nabla F(w_i)\Bigr\rangle .
\]
Applying Cauchy–Schwarz,
\[
\ge \frac{1}{\sqrt{\widehat v_{t,\max}^{\;{\rm corr}}}+\epsilon}\,
\big\|\nabla F(w_t)\big\|^{2}\!\sum_{i=1}^{t}\alpha_{t,i}
= \frac{\|\nabla F(w_t)\|^{2}}{\sqrt{\widehat v_{t,\max}^{\;{\rm corr}}}+\epsilon},
\tag{9}
\]
where $\widehat v_{t,\max}^{\;{\rm corr}}:=\max_{i}\widehat v_{t,i}^{\;{\rm corr}}$.
Since every coordinate of $\widehat v_{t}^{\;{\rm corr}}$ is bounded below by Lemma~\ref{lem:denom},
\[
\sqrt{\widehat v_{t,\max}^{\;{\rm corr}}}+\epsilon\;\ge\;\sqrt{(1-\beta_2)}\,G . \tag{10}
\]

\noindent\textbf{[Step 6] Upper bound the quadratic term.}
From Lemma~\ref{lem:denom} and Lemma~\ref{lem:bias},
\[
\big\|\frac{\widehat m_t}{\sqrt{\widehat v_t^{\;{\rm corr}}}+\epsilon}\big\|^{2}
\le \frac{\|m_t\|^{2}}{(1-\beta_2)G^{2}} .
\]
Using $\|m_t\|\le\frac{1-\beta_1^{t}}{1-\beta_1}\,G$ (derived from the convex combination of bounded $g_i$'s),
\[
\big\|\frac{\widehat m_t}{\sqrt{\widehat v_t^{\;{\rm corr}}}+\epsilon}\big\|^{2}
\le \frac{G^{2}}{(1-\beta_2)(1-\beta_1)^{2}} . \tag{11}
\]

\noindent\textbf{[Step 7] Combine (8)–(11).}
Plugging (9)–(11) into (8) yields
\[
\mathbb{E}\!\bigl[F(w_{t+1})\mid\mathcal{F}_{t-1}\bigr]
\le F(w_t) -\frac{\eta_t}{\sqrt{(1-\beta_2)}\,G}\,\|\nabla F(w_t)\|^{2}
+ \frac{L\eta_t^{2}}{2}\,\frac{G^{2}}{(1-\beta_2)(1-\beta_1)^{2}} . \tag{12}
\]

\noindent\textbf{[Step 8] Take full expectation and sum.}
Taking total expectation and summing (12) for $t=1,\dots,T$,
\[
\mathbb{E}[F(w_{T+1})] \le F(w_1) -\frac{1}{\sqrt{(1-\beta_2)}\,G}
\sum_{t=1}^{T}\eta_t\mathbb{E}\!\bigl[\|\nabla F(w_t)\|^{2}\bigr]
+ \frac{L G^{2}}{2(1-\beta_2)(1-\beta_1)^{2}}
\sum_{t=1}^{T}\eta_t^{2}. \tag{13}
\]

\noindent\textbf{[Step 9] Plug the schedule $\eta_t=\eta/\sqrt{t}$.}
We have
\[
\sum_{t=1}^{T}\eta_t = \eta\sum_{t=1}^{T}\frac{1}{\sqrt{t}}
\ge \eta\int_{1}^{T+1}u^{-1/2}\,du
= 2\eta\bigl(\sqrt{T+1}-1\bigr)
\ge \eta\sqrt{T}. \tag{14}
\]
Similarly,
\[
\sum_{t=1}^{T}\eta_t^{2}= \eta^{2}\sum_{t=1}^{T}\frac{1}{t}
\le \eta^{2}\bigl(1+\log T\bigr). \tag{15}
\]

\noindent\textbf{[Step 10] Rearrange (13).}
Discard the non‑negative term $\mathbb{E}[F(w_{T+1})]$ and use (14)–(15):
\[
\frac{\eta}{\sqrt{(1-\beta_2)}\,G}\,\sqrt{T}\;
\frac{1}{T}\sum_{t=1}^{T}\mathbb{E}\!\bigl[\|\nabla F(w_t)\|^{2}\bigr]
\le F(w_1)-F^\star
+ \frac{L G^{2}\eta^{2}\bigl(1+\log T\bigr)}{2(1-\beta_2)(1-\beta_1)^{2}} .
\]
Dividing by $\eta\sqrt{T}/\bigl(\sqrt{(1-\beta_2)}\,G\bigr)$ gives
\[
\frac{1}{T}\sum_{t=1}^{T}\mathbb{E}\!\bigl[\|\nabla F(w_t)\|^{2}\bigr]
\le
\frac{ \sqrt{(1-\beta_2)}\,G\bigl(F(w_1)-F^\star\bigr)}{\eta\sqrt{T}}
+\frac{L G^{3}\eta\bigl(1+\log T\bigr)}{2(1-\beta_2)^{3/2}(1-\beta_1)^{2}\sqrt{T}} .
\tag{16}
\]

\noindent\textbf{[Step 11] Simplify constants.}
Define
\[
C:=\frac{2(1-\beta_1)}{(1-\beta_2)(1-\beta_1/\sqrt{\beta_2})},
\]
which upper bounds the factor $\displaystyle\frac{G^{3}}{(1-\beta_2)^{3/2}(1-\beta_1)^{2}}$
under the standard condition $\beta_1<\sqrt{\beta_2}$.  Rewriting (16) with this $C$ yields exactly the bound stated in Theorem~\ref{thm:main}.

\noindent\textbf{[Step 12] Obtain the minimum‑gradient guarantee.}
Since the average of non‑negative numbers bounds their minimum,
\[
\min_{1\le t\le T}\mathbb{E}\!\bigl[\|\nabla F(w_t)\|^{2}\bigr]
\le \frac{1}{T}\sum_{t=1}^{T}\mathbb{E}\!\bigl[\|\nabla F(w_t)\|^{2}\bigr]
= O\!\Bigl(\frac{\log T}{\sqrt{T}}\Bigr).
\]
This completes the proof. \qed

%%%%%%%%%%%%%%%%%%%%%%%%%%%%%%%%%%%%%%%%%%%%%%%%%%%%%%%%%%%%
\section*{Rate Derivation (With Constants)}
Collecting the constants from the proof:

\[
\boxed{
\begin{aligned}
A &:= 2L\bigl(F(w_1)-F^\star\bigr),\\[2pt]
B &:= C\,\eta\,G\;(1+\log T),\\[2pt]
\text{Bound} &:= \frac{A}{\eta\sqrt{T}}+\frac{B}{\sqrt{T}}.
\end{aligned}}
\]

Thus the dominant term is $O\!\bigl(\log T/\sqrt{T}\bigr)$.

%%%%%%%%%%%%%%%%%%%%%%%%%%%%%%%%%%%%%%%%%%%%%%%%%%%%%%%%%%%%
\section*{Special Cases}
\begin{enumerate}[leftmargin=*]
\item \textbf{Deterministic gradients ($\sigma=0$).} The variance term disappears and the rate improves to $O(1/\sqrt{T})$.
\item \textbf{Full‑batch Adam ($\widehat v_t=v_t$).} The monotonicity guarantee is lost; the proof above no longer holds because inequality (10) may be violated, leading to possible divergence (cf. Reddi et al., 2018).
\item \textbf{Setting $\beta_1=0$.} AMSGrad becomes a pure RMSProp‑type method with
\[
C=\frac{2}{1-\beta_2},
\]
and the bound simplifies accordingly.
\end{enumerate}

%%%%%%%%%%%%%%%%%%%%%%%%%%%%%%%%%%%%%%%%%%%%%%%%%%%%%%%%%%%%
\section*{Discussion}
The proof shows that the only place where the max‑operator $\widehat v_t=\max(\widehat v_{t-1},v_t)$ is used is to obtain the uniform lower bound (10) on the adaptive denominator.  This bound prevents the effective step size from growing arbitrarily large and is the key to establishing the $O(\log T/\sqrt{T})$ guarantee for non‑convex smooth objectives.  The result matches the best known rates for stochastic first‑order methods under the same assumptions, while offering the practical benefit of per‑coordinate adaptive learning rates.

\end{document}